\section{Bodegas y Centros de Distribucion}

\subsection{Picking rapido y asignacion de espacio}

\begin{materia}{volumen optimo para picking rapido}
Sea $I$ el conjunto de SKU. Para cada SKU $i\in I$:
\begin{itemize}
  \item $f_i$ es el flujo volumetrico mensual en $m^3/\text{mes}$.
  \item $v_i$ es el volumen asignado al SKU $i$.
  \item $V$ es el volumen total disponible.
  \item $c_r$ es el costo por reposicion.
\end{itemize}
El modelo clasico es:
\[
  \min_{v_i\ge 0}\; c_r\sum_{i\in I}\frac{f_i}{v_i}
  \quad \text{s.a.}\quad
  \sum_{i\in I}v_i=V.
\]
La condicion de optimalidad entrega:
\[
  v_i^\star=V\frac{\sqrt{f_i}}{\sum_{j\in I}\sqrt{f_j}}.
\]
Si los costos de reposicion son distintos:
\[
  \min \sum_i c_{r,i}\frac{f_i}{v_i},
  \qquad
  v_i^\star=V\frac{\sqrt{c_{r,i}f_i}}{\sum_j\sqrt{c_{r,j}f_j}}.
\]
\end{materia}

\begin{enunciado}{Examen 2016, centro de distribucion}
En otro centro de distribucion se tiene un area de Picking rapido cuyo espacio total disponible es de $1500\ \text{m}^3$, en la cual se busca disponer tres SKU: X, Y, Z y W.

\begin{center}
\begin{tabular}{c|ccc}
\toprule
SKU & UNIDADES/MES & UNIDADES/CAJA & M3/CAJA\\
\midrule
X & 3000 & 20 & 3\\
Y & 3600 & 24 & 2\\
Z & 2040 & 36 & 1.5\\
W & 1000 & 50 & 5\\
\bottomrule
\end{tabular}
\end{center}

\begin{enumerate}[label=\alph*)]
  \item Plantee el problema de optimizacion que minimiza el costo total del area de picking rapido. Suponga que el costo de \emph{restock} es $c_r$.
  \item Con los datos de la tabla anterior calcule el espacio optimo y la frecuencia de \emph{restock} para cada SKU.
  \item Calcule el espacio asignado y el \emph{restock} para cada SKU considerando las siguientes configuraciones: igual volumen para los SKU y con igual frecuencia de \emph{restock}.
\end{enumerate}
\end{enunciado}

\begin{pautacompleta}{pauta paso a paso 2016}
Primero convierte cada SKU a flujo volumetrico:
\[
  f_i=\frac{\text{unidades/mes}_i}{\text{unidades/caja}_i}\cdot m_i^3/\text{caja}.
\]
Con los datos de la tabla:
\[
  f_X=\frac{3000}{20}\cdot 3=450,\qquad
  f_Y=\frac{3600}{24}\cdot 2=300,
\]
\[
  f_Z=\frac{2040}{36}\cdot 1.5=85,\qquad
  f_W=\frac{1000}{50}\cdot 5=100.
\]
El modelo que minimiza el costo total de reposiciones es:
\[
  \min_{v_i\geq 0}\; c_r\sum_{i\in I}\frac{f_i}{v_i}
  \quad\text{s.a.}\quad
  \sum_{i\in I} v_i=1500,
\]
donde $I=\{X,Y,Z,W\}$, $f_i$ es el flujo volumetrico mensual del SKU $i$ y $v_i$ es el volumen asignado al SKU $i$.

La condicion de optimalidad entrega:
\[
  v_i^\star=1500\frac{\sqrt{f_i}}{\sqrt{450}+\sqrt{300}+\sqrt{85}+\sqrt{100}}.
\]
La suma de raices es:
\[
  \sqrt{450}+\sqrt{300}+\sqrt{85}+\sqrt{100}=57.754.
\]
Los volumenes optimos y frecuencias son:
\[
\begin{array}{c|cc}
\text{SKU} & v_i^\star & \text{restock } f_i/v_i\\
\hline
X&550.96&0.82\\
Y&449.86&0.67\\
Z&239.46&0.35\\
W&259.73&0.39
\end{array}
\]

Para igual volumen:
\[
  v_i=\frac{1500}{4}=375\ \text{m}^3\quad \forall i.
\]
Las frecuencias de reposicion son:
\[
\begin{array}{c|cc}
\text{SKU} & v_i & f_i/v_i\\
\hline
X&375&1.20\\
Y&375&0.80\\
Z&375&0.23\\
W&375&0.27
\end{array}
\]

Para igual frecuencia, se exige $f_i/v_i=k$ para todo $i$, por lo tanto:
\[
  v_i=V\frac{f_i}{\sum_j f_j}.
\]
Como $\sum_j f_j=935$, se obtiene:
\[
\begin{array}{c|cc}
\text{SKU} & v_i & f_i/v_i\\
\hline
X&1500(450/935)=721.93&0.623\\
Y&1500(300/935)=481.28&0.623\\
Z&1500(85/935)=136.36&0.623\\
W&1500(100/935)=160.43&0.623
\end{array}
\]
\end{pautacompleta}

\begin{enunciado}{Examen 2021, picking rapido}
Su empresa ha construido una zona de \emph{picking} rapido de $20\ \text{m}^3$ y se le pide a usted ubicar los SKU que se detallan en la tabla a continuacion.

\begin{center}
\begin{tabular}{c|cc}
\toprule
SKU & cajas/mes & $\text{m}^3$/caja\\
\midrule
A & 40 & 2\\
B & 20 & 1.5\\
C & 45 & 1\\
\bottomrule
\end{tabular}
\end{center}

Si el costo de reposicion es de \$3 por reposicion para cada SKU y usando una politica de volumen optimo y otra de igual tiempo, determine:

\begin{enumerate}[label=\alph*)]
  \item Para cada politica, cuanto espacio asigna a cada uno y cual es el costo de cada politica.
  \item Si los costos de cada reposicion son distintos para cada SKU, siendo \$3 para SKU A, \$2 para SKU B y \$3 para SKU C, plantee un modelo que permita determinar la asignacion optima para cada SKU.
\end{enumerate}
\end{enunciado}

\begin{pautacompleta}{pauta paso a paso 2021}
Se calculan los flujos:
\[
  f_A=40(2)=80,\qquad f_B=20(1.5)=30,\qquad f_C=45(1)=45.
\]
Con volumen optimo:
\[
  v_i^\star=V\frac{\sqrt{f_i}}{\sum_j \sqrt{f_j}}.
\]
Con $V=20$:
\[
  \sum_j\sqrt{f_j}=\sqrt{80}+\sqrt{30}+\sqrt{45}=21.131.
\]
Por lo tanto:
\[
  v_A^\star=20\frac{\sqrt{80}}{21.131}=8.47,\quad
  v_B^\star=20\frac{\sqrt{30}}{21.131}=5.18,\quad
  v_C^\star=20\frac{\sqrt{45}}{21.131}=6.35.
\]
Las reposiciones mensuales son:
\[
  r_A=\frac{80}{8.47}=9.45,\quad
  r_B=\frac{30}{5.18}=5.79,\quad
  r_C=\frac{45}{6.35}=7.09.
\]
El costo mensual de reposicion con volumen optimo es:
\[
  CT_{\text{opt}}=3(r_A+r_B+r_C)=3(22.33)=66.99.
\]

En igual tiempo se igualan las frecuencias:
\[
  \frac{f_A}{v_A}=\frac{f_B}{v_B}=\frac{f_C}{v_C}.
\]
Entonces se asigna proporcional a flujo:
\[
  v_i=V\frac{f_i}{\sum_j f_j},\qquad \sum_j f_j=155.
\]
La tabla queda:
\[
\begin{array}{c|ccc}
\text{SKU} & f_i & v_i & f_i/v_i\\
\hline
A&80&20(80/155)=10.32&7.75\\
B&30&20(30/155)=3.87&7.75\\
C&45&20(45/155)=5.81&7.75
\end{array}
\]
El costo mensual es:
\[
  CT_{\text{igual tiempo}}=3(7.75+7.75+7.75)=69.75.
\]

\textbf{Nota de consistencia.} La pauta escrita usa $V=25$ en los calculos numericos y obtiene $v_A=10.58$, $v_B=6.48$, $v_C=7.94$, costo $53.575$, y para igual tiempo volumenes $12.90$, $4.84$, $7.26$ con costo $55.8$. El enunciado transcrito dice $20\ \text{m}^3$; por eso aqui se recalcula con $V=20$. La proporcion y el algoritmo son los mismos.

Si los costos de reposicion son distintos, el modelo correcto es:
\[
  \min_{v_A,v_B,v_C\geq 0}\;
  3\frac{80}{v_A}+2\frac{30}{v_B}+3\frac{45}{v_C}
\]
sujeto a:
\[
  v_A+v_B+v_C\leq 20,
  \qquad
  v_A,v_B,v_C\geq 0.
\]
En forma general:
\[
  \min_{v_i\geq 0}\sum_{i\in I} c_{r,i}\frac{f_i}{v_i}
  \quad\text{s.a.}\quad
  \sum_{i\in I}v_i\leq V.
\]
\end{pautacompleta}

\subsection{Profundidad optima de pallets}

\begin{materia}{profundidad}
Para SKU $i$:
\[
  d_i^\star=\sqrt{\frac{a q_i}{2z_i}},
\]
donde $a$ es ancho de pasillo medido en pallets, $q_i$ es cantidad de pallets demandados o recibidos, y $z_i$ es altura maxima de apilado. Si hay profundidad unica, se agregan los SKU y se calcula una profundidad comun segun el uso de espacio equivalente.
\end{materia}

\begin{enunciado}{Examen 2022, profundidad de pallets}
Usted debe determinar la profundidad optima de los pallets en una bodega que tiene unos pasillos de 3.5 metros de ancho. Para ello considera una demanda constante y que los pallets tienen una medida de $1\times 1$ metros. Usted recopila la informacion de los 4 SKU que maneja.

\begin{center}
\begin{tabular}{c|cc}
\toprule
SKU & \# Pallets Demandados & Max. Altura de Apilado\\
\midrule
A & 24 & 3\\
B & 20 & 4\\
C & 10 & 2\\
D & 12 & 2\\
\bottomrule
\end{tabular}
\end{center}

\begin{enumerate}[label=\roman*.]
  \item Determine la profundidad optima para cada SKU individualmente.
  \item Determine la profundidad optima para todos en conjunto en el cual hay un pasillo y no se comparte. Muestre sus calculos.
  \item Si ahora puede tener 2 profundidades distintas y comparten el pasillo, cuales serian estas profundidades? Mejora el uso de espacio y por que?
\end{enumerate}
\end{enunciado}

\begin{pautacompleta}{pauta paso a paso 2022}
Como $a=3.5$:
\[
  d_A=\sqrt{\frac{3.5(24)}{2(3)}}=3.74,\quad
  d_B=\sqrt{\frac{3.5(20)}{2(4)}}=2.96,
\]
\[
  d_C=\sqrt{\frac{3.5(10)}{2(2)}}=2.96,\quad
  d_D=\sqrt{\frac{3.5(12)}{2(2)}}=3.24.
\]
La tabla de la pauta es:
\[
\begin{array}{c|cccc}
\text{SKU} & \#\text{ pallets} & \text{altura max.} & \text{prof.} & \text{efectiva}\\
\hline
A&24&3&3.7417&4\\
B&20&4&2.9580&3\\
C&10&2&2.9580&3\\
D&12&2&3.2404&4
\end{array}
\]

Para una profundidad comun sin compartir pasillo, la pauta agrega el espacio equivalente:
\[
\begin{array}{c|ccc}
\text{SKU} & \#\text{ pallets} & \text{altura max.} & \text{espacio}\\
\hline
A&24&3&8\\
B&20&4&5\\
C&10&2&5\\
D&12&2&6\\
\hline
\multicolumn{3}{r}{\text{total}}&24
\end{array}
\]
La profundidad continua reportada es:
\[
  d=4.5826.
\]
Por lo tanto, la profundidad operativa comun es:
\[
  d=5.
\]

Para dos profundidades, se agrupan los SKU con profundidades individuales parecidas:
\[
  \{A,D\}\quad \text{y}\quad \{B,C\}.
\]
La pauta resume:
\[
\begin{array}{c|ccc|c|ccc}
\text{SKU} & \# & h & \text{esp.} & \text{SKU} & \# & h & \text{esp.}\\
\hline
A&24&3&8 & B&20&4&5\\
D&12&2&6 & C&10&2&5\\
\hline
\multicolumn{3}{r}{\text{total}}&14 & \multicolumn{3}{r}{\text{total}}&10
\end{array}
\]
Esto entrega una profundidad de 4 para el grupo $\{A,D\}$ y una profundidad de 3 para el grupo $\{B,C\}$. Mejora el uso de espacio porque evita obligar a los SKU B y C, que naturalmente requieren profundidad cercana a 3, a operar con profundidad comun 5.
\end{pautacompleta}

\subsection{Tamaño de bodega por flujo}

\begin{materia}{flujo, area y trabajadores}
Usa la relacion:
\[
  Q=A\cdot v,
\]
donde $Q$ es flujo anual o semanal, $A$ es inventario/area disponible expresada en pallets, y $v$ es rotacion. Si $T$ es numero de trabajadores, muchas preguntas modelan $v$ como proporcional a $T$.
\end{materia}

\begin{enunciado}{Examen 2021, tamaño minimo de bodega}
Usted se encuentra determinando el tamaño minimo de bodega que debe adquirir. Para ello usted determina que la demanda anual es de 60.000 pallets y los operadores de montacargas trabajan turnos de 8 hrs. por 250 dias al año, con un sueldo mensual de \$15 mil. El tiempo medio desde la recepcion a bodega y a despacho es de 20 minutos. Si el costo anual por $\text{mt}^2$ de terreno e infraestructura para bodega es de \$1.000 el $\text{mt}^2$ y cada pallet utiliza $4\ \text{mt}^2$ de piso en la bodega, ya que pueden ser apilados, caracterice la funcion de costos totales anuales de la bodega. Determine la cantidad de trabajadores y el tamaño optimo de bodega. Muestre todos sus calculos.
\end{enunciado}

\begin{pautacompleta}{pauta paso a paso 2021}
Minutos anuales por operador:
\[
  60\cdot 8\cdot 250=120{,}000.
\]
Como cada pallet requiere 20 minutos, un operador mueve:
\[
  \frac{120{,}000}{20}=6{,}000\;\text{pallets/año}.
\]
Si la demanda es $60{,}000$, la rotacion por operador se usa como:
\[
  \frac{60{,}000}{6{,}000}=10\;\text{rotaciones/año}.
\]
Con $T$ trabajadores:
\[
  v=10T,\qquad A=\frac{60{,}000}{10T}.
\]
El costo de espacio por pallet es:
\[
  1000\ \frac{\$}{m^2}\cdot 4\ \frac{m^2}{\text{pallet}}
  =4000\ \frac{\$}{\text{pallet}}.
\]
La pauta escrita transforma este costo como 250 \$/pallet y plantea:
\[
  CT(T)=250\frac{60{,}000}{10T}+15{,}000T.
\]
Siguiendo esa pauta:
\[
  CT'(T)=-\frac{1{,}500{,}000}{T^2}+15{,}000=0
  \Rightarrow T=10.
\]
Entonces:
\[
  A=\frac{60{,}000}{10(10)}=600\;\text{pallets},
  \qquad
  \text{area}=600(4)=2400\,m^2.
\]

\textbf{Nota de consistencia.} El enunciado dice costo anual de \$1.000 por $m^2$ y cada pallet usa $4m^2$, lo que aritmeticamente da \$4.000 por pallet, no \$250. La pauta antigua usa \$250/pallet para llegar a $T=10$ y $2400m^2$. Para estudiar la mecanica del examen, lo importante es plantear $CT(T)$ y derivar; si se usa literalmente el costo del enunciado, cambia el optimo numerico.
\end{pautacompleta}
